%! Bias in Samples and Surveys — Unit 1, Lesson 1.4 — Homework
% PILOT SHAPE (2026-09-06): modeled on AP Statistics lesson 1.4 minus its AP Practice page.
% This is the lesson's SCORED individual practice and the LAST component in the packet.
% Students START IT IN CLASS in the last ten minutes of the period, alone, while the
% teacher circulates for the second formative read; they finish it at home. Packet pages are
% the default; a Desmos activity may replace them on a given day (decided at Close & Assign).
% THIRD CONTEXT: the notes teach on Harbor Point pool, the Guided Practice runs on Cedar Ridge
% Apartments, and these items run on Millbrook Public Library — recall is not the skill.
% TWO PAGES MAX (user decision 2026-09-06): two boxes — items 1-2 on page 1, items 3-6 and
% the spiralbox on page 2 — so no item breaks across a page.
% Item 4 is the deliberate CONTRAST PAIR (Plan A, four weeks at the desk / Plan B, the phone
% sample: bigger vs. closer to the truth) and then the target misconception head-on — what
% four weeks of the same flawed collection would report. Item 1 is the spiral to 1.2/1.3
% (population, sample, who chose). Item 6 is AFDA.DA.2g: a random draw with a 65% response
% rate can still be biased; name the fix.
% THE DATA — Millbrook Public Library: 1,200 card holders. Front desk: 72/90 = 80% -> 960.
%   Phone sample of 120, every one reached: 54/120 = 45% -> 540. Gap 35 points, 420 holders.
%   Question order: 68% vs 47% = 21 points. Response rate 78/120 = 65%; 42 never picked up.
% PAGE LOCKSTEP: the key is generated from this file; every work block is byte-identical.
\documentclass[12pt]{article}
\usepackage{saar-article}
\usepackage{saar-boxes}
\newcommand{\TallMath}[1]{$\displaystyle #1\rule[-1.4em]{0pt}{3.2em}$}
% Word bank strip for every fill-in cluster (user request 2026-09-06). Defined per document;
% the key inherits it and prints the same strip, so heights match.
\newcommand{\wordbank}[1]{\par\vspace{2pt}\noindent\fcolorbox{goldacc}{goldbg}{\parbox{\dimexpr\linewidth-2\fboxsep-2\fboxrule\relax}{\small\textbf{Word bank:} #1}}\par\vspace{4pt}}

\begin{document}

\pageheader{Unit 1, Lesson 1.4}{Homework}
% NAMESTRIP: no \namedateperiod here. The student writes their name once, on the
% cover this component is stapled behind.

\vspace{0.12in}

% Authored identically in homework/ and homework_key/.
\begin{remindbox}
\small
\textbf{This is your graded homework.} It is scored out of 12 and due at the start of the
first class after two study halls. You will start it in the last minutes of the period, so ask about anything that stalls
you before you leave, then finish it on your own.
\end{remindbox}

\vspace{0.08in}

\boxguard[20]
\begin{scenariobox}[Millbrook Public Library --- Should Weekend Hours Be Longer?]{cerulean}
Millbrook Public Library has $1{,}200$ card holders. The director wants to know what share of
them want longer weekend hours before she asks the town for the money.

\vspace{2pt}
Her first attempt: a stack of paper surveys left on the front desk for one week, with a box
beside it. $90$ surveys come back, and $72$ of those say they want longer weekend hours.
\end{scenariobox}

\vspace{0.08in}

\boxguard
\begin{notesbox}{Two Studies, Two Answers}
Every item below is about the same question --- the studies differ only in \emph{who got
asked} and \emph{who answered}.
\wordbank{all 1,200 card holders \;\textbullet\; the 90 who filled one out \;\textbullet\; convenience \;\textbullet\; undercoverage \;\textbullet\; 35 \;\textbullet\; 420 \;\textbullet\; A \;\textbullet\; B \;\textbullet\; about 80\% \;\textbullet\; about 45\% \;\textbullet\; acquiescence \;\textbullet\; demand \;\textbullet\; social desirability \;\textbullet\; neutral responding \;\textbullet\; 21 \;\textbullet\; question order \;\textbullet\; 65\%}
\begin{enumerate}[leftmargin=1.6em, itemsep=9pt, topsep=3pt]

  \item The front-desk study.
        \par\vspace{4pt}
        Population: \blank{4.2cm} \hfill Sample: \blank{5cm}
        \par\vspace{8pt}
        It was a \blank{2.6cm} sample, and a frame that leaves out part of the population
        causes \blank{3cm}. Name one kind of card holder who had almost no chance of being
        included.
        \par\writeline

  \item Of the $90$ surveys returned, $72$ said they want longer weekend hours.

        \textbf{(a)} What percent of the returned surveys is that?

        \begin{work}
          \dfrac{72}{90} &= 0.80 = 80\%
        \end{work}

        \textbf{(b)} The director applies that percent to all $1{,}200$ card holders. How many
        does she predict want longer weekend hours?

        \begin{work}
          0.80 \times 1200 &= 960 \text{ card holders}
        \end{work}

\end{enumerate}
\end{notesbox}

\boxguard[30]
\begin{notesbox}{Two Studies, Two Answers, continued}
\begin{enumerate}[leftmargin=1.6em, itemsep=2pt, topsep=2pt, start=3]

  \item A statistician talks her into doing it properly: a random sample of $120$ card holders
        from the full list, reached by phone until every one answers. $54$ want longer hours.

        \textbf{(a)} What percent is that, and what does it predict for all $1{,}200$?

        \begin{work}
          \dfrac{54}{120} &= 0.45 = 45\%, \text{ so } 0.45 \times 1200 = 540 \text{ card holders}
        \end{work}

        \textbf{(b)} The front-desk study was too high by \blank{1.2cm} points --- about
        \blank{1.4cm} card holders.

  \item Before the phone survey, the director had said: ``The front-desk number is fine, we
        just need more of them --- leave the surveys out for four weeks and get $360$ back.''
        Call that Plan A; the phone sample is Plan B.
        \par\vspace{2pt}
        Bigger: Plan \blank{0.8cm}. \quad Closer to the truth: Plan \blank{0.8cm}. \quad
        Plan A's $360$ would report \blank{2cm} --- why?
        \par\writelines{2}

  \item Name the kind of response bias in each situation.
        \nopagebreak

        \begin{center}
        \renewcommand{\arraystretch}{1.15}
        \begin{tabularx}{\linewidth}{@{} X p{3.6cm} @{}}
        \toprule
        \textbf{The situation} & \textbf{Kind} \\
        \midrule
        A card holder agrees with all ten statements, two of which conflict. & \blank{3.3cm} \\
        The survey is headed \emph{``Help us win longer weekend hours!''} & \blank{3.3cm} \\
        The story-hour librarian hands out the survey and waits for it. & \blank{3.3cm} \\
        A card holder circles the middle box, ``no opinion,'' on every question. & \blank{3.3cm} \\
        \bottomrule
        \end{tabularx}
        \end{center}

        Asked \emph{``Have you ever found the library closed?''} first, $68\%$ support longer
        hours; asked \emph{``Should the town raise library funding?''} first, $47\%$ do. Gap:
        \blank{1cm} points --- \blank{2.9cm} bias.

  \item Same phone draw, but only $78$ of the $120$ ever picked up. Percent answered:
        \blank{1.4cm} \quad The draw was random and every count is correct --- explain why
        the result could still be biased, and name the fix.
        \par\writelines{2}

\end{enumerate}
\end{notesbox}

\boxguard[5]   % the two-line spiralbox needs about five baselines; a larger guard pushes it to a third page
\begin{spiralbox}
\textbf{Next lesson: Observational Studies.} Every study in this unit so far has only
\emph{watched} --- nobody changed anything. What can a study that only watches prove?
\end{spiralbox}

\end{document}
